% ============================================================================
% Title:        Summary
% Author:       Nathaniel Thomas
% Affiliation:  Texas A&M University, Department of Nuclear Engineering
% Email:        nathaniel@tamu.edu
% Date Created: 3/7/2026
% Version:      1.0
%
% Description:
%   A .tex source file containing the summary of the FLiNaK purifier design 
%   document.
%
% =============================================================================

\section{Summary}
This document outlines the design, construction, and operation of the
Texas A\&M FLiNaK purifier. This purifier utilizes
hydrofluorination to reduce metallic impurites, precipitate them, and filter
them out, while converting oxygen and similar non-metallic contaminants to
gaseous form and sparging them out of the molten salt. This process results in
very high purity salt that can be used for corrosion, characterization, and
other experiments. \par

Included in this document is an introductory overview of the system, where the
operating principles and concepts are explained. It also goes over previous
efforts to purify flouride salts and how those efforts succeeded, as well as
how they failed. Following that, there is a section on the detailed design and
construction of the purifier. The design decisions behind the construction, as
well as how the purifier is manufactured is covered. This includes
drawings for custom-fabricated parts, as well as tubing, valves, electronics,
etc used to control gas flow and transfer purified salt from the purifier. \par

In addition to the operating prinicples and design, there is also a section
detailing the programming used to control the purifier. Much of the control
system is made using LabView, in combination with a LabView Digital to Analog
converter (DAC), which enables the control and monitoring of the system.
The program reads parameters such as pressure, temperature, and HF
concentration, and will raise alarms, or at worst, stop all gas flow if
any unusual operating parameters are detected. \par

Finally, a section is dedicated to the safe and correct operating procedure for
the system. It includes a Failure Mode and Effects Analysis (FMEA) which
addresses potential failures that may occur during operation, and what actions
or engineering barriers exists to mitigate or eliminate risks of each potential
failure. An appendix then contains supplementary materials that is referred
to throughout the document. \par

\pagebreak
\tableofcontents
\pagebreak
